\documentclass[11pt]{article}
\usepackage{amsmath,amssymb,amsthm}
\usepackage{booktabs}
\usepackage[margin=1in]{geometry}

\newtheorem{theorem}{Theorem}
\newtheorem{lemma}[theorem]{Lemma}

\title{Problem 682: 5-Smooth Pairs}
\author{Project Euler Solutions}
\date{}

\begin{document}
\maketitle

\section{Problem Statement}

A positive integer is \emph{5-smooth} if its largest prime factor is at most~5. For a positive integer~$a$, let $\Omega(a)$ count its prime factors with multiplicity and $s(a)$ their sum. Define $f(n)$ as the number of ordered pairs $(p,q)$ of 5-smooth numbers with $\Omega(p)=\Omega(q)$ and $s(p)+s(q)=n$.

Given: $f(10)=4$, $f(100)=3629$. Find $f(10^7)\bmod(10^9+7)$.

\section{Mathematical Foundation}

\begin{theorem}[Generating Function]
Every 5-smooth number $m = 2^a\cdot 3^b\cdot 5^c$ satisfies $\Omega(m)=a+b+c$ and $s(m)=2a+3b+5c$. The bivariate generating function is
\[
G(x,y) = \frac{1}{(1-xy^2)(1-xy^3)(1-xy^5)},
\]
where $[x^k y^j]\,G(x,y)$ counts 5-smooth numbers with $\Omega=k$ and $s=j$.
\end{theorem}

\begin{proof}
The factor $1/(1-xy^p)=\sum_{e\ge 0}x^e y^{pe}$ tracks $e$ copies of prime~$p$, contributing $e$ to $\Omega$ and $pe$ to~$s$. The product over $p\in\{2,3,5\}$ generates all triples $(a,b,c)$.
\end{proof}

\begin{lemma}[Polynomial $h_k$]
Define $h_k(y) = [x^k]\,G(x,y) = \sum_{\substack{a+b+c=k\\a,b,c\ge 0}} y^{2a+3b+5c}$. Then $h_k$ has minimum degree $2k$ and maximum degree $5k$.
\end{lemma}

\begin{proof}
The minimum of $2a+3b+5c$ subject to $a+b+c=k$ is $2k$ (set $a=k$). The maximum is $5k$ (set $c=k$).
\end{proof}

\begin{theorem}[Pair Counting]
$\displaystyle f(n) = \sum_{k=0}^{\lfloor n/2\rfloor} [y^n]\,h_k(y)^2$.
\end{theorem}

\begin{proof}
A pair $(p,q)$ with $\Omega(p)=\Omega(q)=k$ and $s(p)+s(q)=n$ corresponds to choosing $s(p)=j$ and $s(q)=n-j$ from the distribution $h_k$. The number of such pairs is $[y^n]\,h_k(y)^2$ by the convolution theorem. Summing over~$k$ yields $f(n)$.
\end{proof}

\begin{lemma}[Verification]
$f(10) = 4$: contributions from $k=1$ ($[y^{10}]\,h_1^2=1$) and $k=2$ ($[y^{10}]\,h_2^2=3$).
\end{lemma}

\begin{proof}
$h_1(y) = y^2+y^3+y^5$, so $h_1^2 = y^4+2y^5+y^6+2y^7+2y^8+y^{10}$ and $[y^{10}]=1$. For $h_2$, the pairs $(s_1,s_2)$ summing to~10 with $s_i \in \{4,5,6,7,8,10\}$ are $(4,6),(5,5),(6,4)$, giving $[y^{10}]=3$. For $k\ge 3$, $\min(s(p)+s(q))=4k\ge 12>10$.
\end{proof}

\section{Algorithm}

\begin{enumerate}
  \item For each $k = 0,1,\ldots,\lfloor n/2\rfloor$: build $h_k$ as a polynomial, square it, extract $[y^n]$.
  \item Optimization: combine into a single convolution using NTT modulo $10^9+7$.
\end{enumerate}

\section{Complexity Analysis}

\begin{itemize}
  \item \textbf{Time:} $O(n\log n)$ with NTT-based combined convolution; $O(n^2\log n)$ naive per-$k$.
  \item \textbf{Space:} $O(n)$.
\end{itemize}

\section{Answer}

\[
\boxed{290872710}
\]

\end{document}
